\section{Project Maven：检测模型只是 Decision Chain 的入口}

Project Maven 的课堂案例从大范围传感器数据中的目标检测开始。这里不展开攻击操作，而聚焦系统问题：当感兴趣对象只占巨大视野的极小部分，人工逐帧搜索速度不足；computer vision 可以缩短 observe 阶段，却不能自行完成身份确认、法律审核、资源分配、行动授权和结果复盘。一个高准确率 detector 如果没有进入受控 workflow，只会生成更多无人处理的 alert。

\subsection{OODA loop 与 latency budget}

\term{OODA} 是 Observe、Orient、Decide、Act。Observe 收集传感器和环境状态，Orient 将观测放入任务、历史和规则语境，Decide 由有权主体选择行动，Act 执行并产生新结果。循环速度重要，但每个阶段都需要准确性和权限控制。Latency budget 应分配到整个 loop，而不是只优化 inference milliseconds。

\begin{figure}[H]
\centering
\includegraphics[width=0.94\textwidth]{images/08-ooda-decision-chain.png}
\caption{Project Maven 的系统视角：检测加速 Observe，价值来自完整 OODA decision chain。}
\figsource{依据字幕 21:27--27:00 重绘，`images/08-ooda-decision-chain.png`。}
\end{figure}

\begin{knowledgebox}{读图：每条箭头都需要 provenance 与 authority}
Detector 输出应保留传感器、时间、模型版本和 confidence；Orient 阶段需要关联其他情报与不确定性；Decide 阶段必须显示谁有权限、哪些规则适用；Act 后要记录结果与副作用。没有这些链路，模型误报无法追溯，人工判断也无法用于下一轮改进。
\end{knowledgebox}

\begin{importantbox}{术语消化：Decision chain 与 Ontology}
Decision chain 是从观测到行动再到反馈的完整组织流程。Ontology 在这里不是哲学名词，而是共享业务/任务模型：有哪些 entity、它们的关系、可执行 action、状态和权限。Ontology 让传感器、模型、人员和 workflow 使用同一语义层，避免每个系统用自己的 ID 和表结构描述同一对象。
\end{importantbox}

\subsection{从发现对象到受控行动}

讲者描述了从“找到大量候选对象”继续向下追问的过程：如何确认身份，怎样选择资源，何时需要法律或政策审核，结果如何回流。这个扩展揭示了 AI deployment 的常见规律：第一个模型解决可见的局部任务，真正阻塞价值的往往是后续组织流程和数据断点。

\begin{warningbox}{高风险场景中的自动化边界}
检测、排序和建议可以提高人的处理能力，不等于授权系统自行作出不可逆决定。权限、human review、双人复核、置信阈值、申诉或纠错机制应根据后果设计。模型 accuracy 不能替代合法性，也不能消除责任主体。
\end{warningbox}

\subsection{本章小结}

Maven 的教学价值不在某个视觉模型，而在把 inference 放入 OODA loop。观察、定向、决策、行动和反馈需要统一语义、权限和审计。这个骨架也可以用于商业企业，只是目标、合法授权和失败后果不同。

